% ============================================================
% Paper B — Quantifying the Pole-Cancellation Barrier
%           in the Guinand-Weil Explicit Formula
% Version: 1.0 (Threshold assembly, 2026-03-29)
% Status:  Draft — awaiting Henry's decisions on title/authorship/references
% ============================================================
\documentclass[12pt]{article}

% --- Packages ---
\usepackage{amsmath,amssymb,amsthm}
\usepackage{enumerate}
\usepackage{booktabs}
\usepackage[margin=1in]{geometry}
\usepackage[colorlinks=true,linkcolor=blue,citecolor=blue,urlcolor=blue]{hyperref}

% --- Theorem environments ---
\newtheorem{theorem}{Theorem}[section]
\newtheorem{lemma}[theorem]{Lemma}
\newtheorem{corollary}[theorem]{Corollary}
\newtheorem{proposition}[theorem]{Proposition}
\newtheorem{definition}[theorem]{Definition}
\theoremstyle{remark}
\newtheorem{remark}[theorem]{Remark}
\newtheorem{observation}[theorem]{Numerical Observation}

% --- PDF metadata ---
\hypersetup{
  pdftitle={Quantifying the Pole-Cancellation Barrier in the Guinand-Weil Explicit Formula},
  pdfauthor={Henry Chan},
  pdfsubject={Analytic Number Theory, Riemann Hypothesis, Explicit Formulas},
  pdfkeywords={Riemann hypothesis, Guinand-Weil explicit formula, pole cancellation, test function, mollifier barrier}
}

% ============================================================
\title{Quantifying the Pole-Cancellation Barrier\\
in the Guinand-Weil Explicit Formula}

% DECISION PENDING: authorship — Henry Chan solo (AI in Acknowledgments)
% vs. "Gemini-Omega and Henry Chan"
\author{Henry Chan\thanks{B2AGI / exist.is. Contact: henry@b2agi.is}}
\date{Preprint --- April 2026}

\begin{document}
\maketitle

% ============================================================
% ABSTRACT
% ============================================================
\begin{abstract}
We study the pole-cancellation problem in the Guinand--Weil explicit formula
for the Riemann zeta function. For a parametrized family of Gaussian test
functions $F_X$ with scale parameter $X$, we construct a linearly corrected
function $G_X^*$ satisfying the pole-cancellation condition
$\Phi_{G_X^*}(1)=0$ and analyze the resulting contamination-to-signal ratio
$R(X,\varepsilon)$ at a hypothetical off-critical zero
$\rho_0 = \tfrac{1}{2}+\varepsilon+i\gamma$. We prove that
$\log R / X \to f(\varepsilon)$ as $X\to\infty$, where
$f(\varepsilon) = \tfrac{5}{16} - \tfrac{\varepsilon}{2} -
\tfrac{\varepsilon^2}{4}$ is a structural exponent that vanishes exactly at
$\varepsilon^* = \tfrac{1}{2}$---the critical line itself. This pole-zero
exponent inequality (Theorem~\ref{thm:wall}) establishes that for any
$\varepsilon \in (0,\tfrac{1}{2})$, pole cancellation in the Gaussian family
induces exponentially growing contamination that overwhelms the off-critical
signal. We formulate a conditional reduction (Theorem~\ref{thm:conditional})
showing that existence of a Schwartz-class test function simultaneously
achieving pole cancellation and off-critical detection would imply the
Riemann Hypothesis. Numerical computations with 50-digit precision confirm
the theoretical predictions. The structural analogy between $f(\varepsilon)=0$
at $\varepsilon=\tfrac{1}{2}$ and the mollifier-length barrier at
$\theta=\tfrac{1}{2}$ in the Levinson--Conrey framework is discussed.
This paper does not prove the Riemann Hypothesis; it measures where
the explicit-formula approach encounters a quantitative obstruction.
\end{abstract}

% ============================================================
% SECTION 1: INTRODUCTION
% ============================================================
\section{Introduction}

\subsection{Context and Motivation}

The Riemann Hypothesis (RH), asserting that all nontrivial zeros of
$\zeta(s)$ satisfy $\mathrm{Re}(s) = \tfrac{1}{2}$, remains one of the
central open problems in mathematics. Among the principal approaches to RH,
the \emph{explicit formula} of Guinand \cite{Guinand1948} and Weil
\cite{Weil1952} occupies a distinguished position: it provides an exact
identity linking the distribution of prime numbers to the location of zeta
zeros, mediated by a choice of test function.

A natural strategy is to engineer a test function $h$ whose Mellin transform
$\Phi_h$ amplifies any hypothetical off-critical zero while suppressing the
dominant contribution from the pole of $\zeta(s)$ at $s=1$. This
pole-cancellation approach has been pursued in various forms, but invariably
encounters a structural obstruction: the correction required to cancel the
pole contaminates the zero-detection signal by an amount that grows
exponentially with the scale parameter.

The present paper makes this obstruction precise. We introduce a scalar
\emph{contamination-to-signal ratio} $R(X,\varepsilon)$ and derive an exact
structural exponent $f(\varepsilon) = \tfrac{5}{16} - \tfrac{\varepsilon}{2}
- \tfrac{\varepsilon^2}{4}$ governing its exponential growth. The exponent
vanishes at $\varepsilon^* = \tfrac{1}{2}$, which is the critical line
itself---the only point where contamination does not overwhelm the signal.

\subsection{Main Results}

Our two principal results are:

\begin{enumerate}[(i)]
\item \textbf{Conditional reduction (Theorem~\ref{thm:conditional}).}
  If there exists a family of Schwartz-class test functions simultaneously
  achieving pole cancellation and exponential off-critical detection, then
  the Riemann Hypothesis holds. This reduces RH to a concrete existence
  problem for test functions.

\item \textbf{The pole-zero exponent inequality (Theorem~\ref{thm:wall}).}
  For the Gaussian test-function family with linear pole cancellation,
  the contamination-to-signal ratio satisfies
  $R(X,\varepsilon) = \Theta(\exp(f(\varepsilon)\cdot X))$ with
  $f(\varepsilon) > 0$ for all $\varepsilon \in (0,\tfrac{1}{2})$.
  Consequently, this canonical family cannot satisfy the conditions of
  Theorem~\ref{thm:conditional}.
\end{enumerate}

\paragraph{What this paper does.}
We quantify a specific structural barrier within the Guinand--Weil explicit
formula framework. The exponent $f(\varepsilon)$ and its vanishing at
$\varepsilon^* = \tfrac{1}{2}$ are, to our knowledge, new.

\paragraph{What this paper does not claim.}
We do not prove the Riemann Hypothesis. Theorem~\ref{thm:conditional}
provides a conditional reduction whose hypotheses remain unverified. We do
not claim that the Gaussian barrier is universal across all test-function
classes. We do not establish a formal equivalence with the mollifier-length
barrier, though we identify a structural analogy.

\subsection{Structural Analogy with Mollifier Theory}

In the Levinson--Conrey approach to zero proportions on the critical line
\cite{Levinson1974,Conrey1989}, mollifiers of length $T^\theta$ encounter
a well-known barrier at $\theta = \tfrac{1}{2}$: for $\theta < \tfrac{1}{2}$,
diagonal terms dominate off-diagonal contributions, yielding positive
proportions (Levinson's $\tfrac{1}{3}$, Conrey's $\tfrac{2}{5}$). Bettin and
Gonek \cite{BettinGonek2016} showed that $\theta \to \infty$ would be
equivalent to RH.

Our framework exhibits a parallel structure: the scale parameter $X$
plays a role analogous to $\theta \cdot \log T$, and the structural exponent
$f(\varepsilon)$ vanishes at $\varepsilon = \tfrac{1}{2}$, mirroring the
mollifier barrier at $\theta = \tfrac{1}{2}$. Both barriers originate from
the pole of $\zeta(s)$ at $s=1$ contaminating the analytic structure of the
detection mechanism. This is an analogy, not a formal equivalence; establishing
a rigorous connection between the two frameworks remains an open problem.

\subsection{Organization}

Section~\ref{sec:prelim} establishes notation and recalls the Guinand--Weil
explicit formula. Section~\ref{sec:main} presents the main results:
Theorem~\ref{thm:conditional} (conditional reduction),
Lemma~\ref{lem:prime_bound} (prime-sum bound),
and Theorem~\ref{thm:wall} (the pole-zero exponent inequality).
Section~\ref{sec:numerical} provides numerical verification with 50-digit
precision. Section~\ref{sec:discussion} discusses connections to mollifier
theory, zero-density estimates, limitations, and open questions.
Appendix~\ref{app:notation} collects notation, and
Appendix~\ref{app:proofs} contains detailed proofs.

% ============================================================
% SECTION 2: PRELIMINARIES
% ============================================================
\section{Preliminaries}\label{sec:prelim}

\subsection{The Guinand--Weil Explicit Formula}

For a test function $h$ in the Schwartz class $\mathcal{S}(\mathbb{R}_{>0})$
with Mellin transform
$\Phi_h(s) = \int_0^\infty h(x)\, x^{s-1}\, dx$,
the Guinand--Weil explicit formula states:
\begin{equation}\label{eq:GW}
\sum_{\rho} \Phi_h(\rho)
= \Phi_h(1) + \Phi_h(0)
  - \sum_{n=2}^{\infty} \frac{\Lambda(n)}{\sqrt{n}}
    \bigl[h(n) + h(1/n)\bigr]
  + W_\infty(h),
\end{equation}
where $\rho$ ranges over the nontrivial zeros of $\zeta(s)$, $\Lambda$ is
the von Mangoldt function, and $W_\infty(h)$ is the archimedean contribution.

\subsection{Gaussian Test Functions}

\begin{definition}\label{def:gaussian}
For a scale parameter $X > 0$, define the Gaussian test function
\begin{equation}\label{eq:FX}
F_X(x) = x^{-1/2}\,
\exp\!\Bigl(-\frac{(\log x - X)^2}{2X}\Bigr).
\end{equation}
\end{definition}

The Mellin transform of $F_X$ evaluates in closed form:
\begin{equation}\label{eq:PhiX}
\Phi_X(s) = \sqrt{2\pi X}\;
\exp\!\Bigl((s - \tfrac{1}{2})\,X
  + \tfrac{1}{2}(s - \tfrac{1}{2})^2\,X\Bigr).
\end{equation}

At $s = 1$:
\begin{equation}\label{eq:Phi1}
|\Phi_X(1)| = \Theta\bigl(X^{1/2} \cdot e^{5X/8}\bigr).
\end{equation}

At $s = \tfrac{1}{2} + \varepsilon + i\gamma$ with $\varepsilon > 0$:
\begin{equation}\label{eq:Phirho}
|\Phi_X(\rho_0)| = \Theta\bigl(X^{1/2} \cdot
  e^{\varepsilon X + \varepsilon^2 X/2}\bigr).
\end{equation}

\subsection{Pole-Cancellation Correction}

To enforce $\Phi_{G_X^*}(1) = 0$, we define:
\begin{equation}\label{eq:correction}
G_X^*(x) = F_X(x) - c_X \cdot F_{X_0}(x),
\qquad c_X = \frac{\Phi_X(1)}{\Phi_{X_0}(1)},
\end{equation}
where $X_0 = \kappa X$ with $\kappa \in (0,1)$ is a reference scale.
By construction, $\Phi_{G_X^*}(1) = \Phi_X(1) - c_X \Phi_{X_0}(1) = 0$.

The correction coefficient grows as:
\begin{equation}\label{eq:cX}
|c_X| = \Theta\bigl(\kappa^{-1/2} \cdot
  e^{5(1-\kappa)X/8}\bigr).
\end{equation}

% ============================================================
% SECTION 3: MAIN RESULTS
% ============================================================
\section{Main Results}\label{sec:main}

\begin{theorem}[Conditional Reduction]\label{thm:conditional}
Let $h \in \mathcal{S}(\mathbb{R}_{>0})$ be a Schwartz-class test function
with Mellin transform $\Phi_h$. Suppose:
\begin{enumerate}[\upshape (C1)]
\item \textbf{Pole cancellation:}
  $\Phi_h(1) = 0$.
\item \textbf{Off-critical detection:}
  There exist $\varepsilon > 0$, $\gamma \in \mathbb{R}$, $C > 0$, and
  $X_0 > 0$ such that for all $X > X_0$,
  \[
  |\Phi_h(\tfrac{1}{2} + \varepsilon + i\gamma)| \geq C \cdot e^{\varepsilon X}.
  \]
\end{enumerate}
Then $\rho_0 = \tfrac{1}{2} + \varepsilon + i\gamma$ is not a nontrivial
zero of $\zeta(s)$.

Consequently, if a family of test functions satisfying \textup{(C1)} and
\textup{(C2)} exists for all $\varepsilon > 0$ and all $\gamma \in \mathbb{R}$,
the Riemann Hypothesis would follow.
\end{theorem}

\begin{proof}[Proof sketch]
Under (C1), the dominant term $\Phi_h(1)$ vanishes from the right-hand side
of \eqref{eq:GW}. By Lemma~\ref{lem:prime_bound}, the remaining right-hand
side is $O(C_h \cdot X)$. Under (C2), the contribution of $\rho_0$ to the
left-hand side is $\Omega(e^{\varepsilon X})$. Since
$e^{\varepsilon X} > C' X$ for all sufficiently large $X$, the explicit
formula \eqref{eq:GW} yields a contradiction. Hence $\rho_0$ cannot be a
zero. Full details appear in Appendix~\ref{app:proofs}.
\end{proof}

\begin{lemma}[Prime-Sum Bound]\label{lem:prime_bound}
Let $h \in \mathcal{S}(\mathbb{R}_{>0})$ satisfy $\Phi_h(1) = 0$.
Then the prime sum on the right-hand side of \eqref{eq:GW} satisfies
\[
\Biggl|\sum_{n=2}^{\infty} \frac{\Lambda(n)}{\sqrt{n}}
  \bigl[h(n) + h(1/n)\bigr]\Biggr|
= O(C_h \cdot X),
\]
where $C_h$ is a constant depending on the Schwartz norms of $h$.
\end{lemma}

\begin{proof}[Proof sketch]
By the prime number theorem, $\psi(x) = x + O(x\exp(-c\sqrt{\log x}))$.
Summation by parts against the rapidly decaying $h$ yields the stated bound.
The Schwartz decay of $h$ ensures absolute convergence and uniformity in $X$.
See Appendix~\ref{app:proofs} for details.
\end{proof}

\begin{corollary}\label{cor:oneway}
If $\rho_0 = \tfrac{1}{2}+\varepsilon+i\gamma$ is a nontrivial zero of
$\zeta(s)$ with $\varepsilon > 0$, then no Schwartz-class test function
can simultaneously satisfy \textup{(C1)} and \textup{(C2)} at $\rho_0$.
\end{corollary}

\begin{remark}
The implication in Theorem~\ref{thm:conditional} is one-directional:
the existence of such test functions would imply RH, but the converse
is not established. Whether RH implies the existence of functions
satisfying (C1)+(C2) is itself an open question.
\end{remark}

\medskip

We now show that the Gaussian family cannot satisfy (C1)+(C2) simultaneously.

\begin{theorem}[The Pole-Zero Exponent Inequality]\label{thm:wall}
Let $F_X$ be the Gaussian test function \eqref{eq:FX} with scale parameter
$X$, and let $G_X^* = F_X - c_X F_{X_0}$ be the linearly corrected function
\eqref{eq:correction} with $X_0 = \kappa X$, $\kappa \in (0,1)$. For any
hypothetical zero $\rho_0 = \tfrac{1}{2}+\varepsilon+i\gamma_0$ with
$\varepsilon > 0$, the contamination-to-signal ratio
\[
R(X,\varepsilon) := \frac{|c_X \cdot \Phi_{X_0}(\rho_0)|}{|\Phi_X(\rho_0)|}
\]
satisfies, for $\kappa = \tfrac{1}{2}$:
\begin{equation}\label{eq:R}
R(X,\varepsilon) = \Theta\bigl(\exp(f(\varepsilon) \cdot X)\bigr),
\end{equation}
where the \emph{structural exponent} is
\begin{equation}\label{eq:feps}
f(\varepsilon) = \frac{5}{16} - \frac{\varepsilon}{2}
  - \frac{\varepsilon^2}{4}.
\end{equation}
In particular:
\begin{enumerate}[\upshape (a)]
\item $f(\varepsilon) > 0$ for all $\varepsilon \in (0, \tfrac{1}{2})$,
  so $R(X,\varepsilon) \to \infty$ as $X \to \infty$.
\item $f(\varepsilon^*) = 0$ at the unique point
  $\varepsilon^* = \tfrac{1}{2}$, which corresponds to the critical line
  $\mathrm{Re}(\rho) = 1$.
\item For general $\kappa \in (0,1)$, the exponent becomes
  $(1-\kappa)\bigl(\tfrac{5}{8} - \varepsilon - \tfrac{\varepsilon^2}{2}\bigr)$,
  which is positive for all $\varepsilon < \tfrac{1}{2}$ and all $\kappa$.
\end{enumerate}

Consequently, the Gaussian family with linear pole cancellation cannot
satisfy condition \textup{(C2)} of Theorem~\ref{thm:conditional} for any
$\varepsilon \in (0,\tfrac{1}{2})$.
\end{theorem}

\begin{proof}
From \eqref{eq:Phi1} and \eqref{eq:Phirho}, at $\kappa = \tfrac{1}{2}$:
\begin{align}
|c_X| &= \frac{|\Phi_X(1)|}{|\Phi_{X/2}(1)|}
  = \Theta\bigl(2^{1/2} \cdot e^{5X/16}\bigr), \label{eq:cX_proof}\\
|\Phi_{X/2}(\rho_0)| &= \Theta\bigl((X/2)^{1/2} \cdot
  e^{\varepsilon X/2 + \varepsilon^2 X/4}\bigr), \label{eq:Phi_ref}\\
|\Phi_X(\rho_0)| &= \Theta\bigl(X^{1/2} \cdot
  e^{\varepsilon X + \varepsilon^2 X/2}\bigr). \label{eq:Phi_main}
\end{align}
Therefore:
\begin{align}
R(X,\varepsilon)
&= \frac{|c_X| \cdot |\Phi_{X/2}(\rho_0)|}{|\Phi_X(\rho_0)|} \notag\\
&= \Theta\Bigl(\exp\bigl(\tfrac{5X}{16}
  + \tfrac{\varepsilon X}{2} + \tfrac{\varepsilon^2 X}{4}
  - \varepsilon X - \tfrac{\varepsilon^2 X}{2}\bigr)\Bigr) \notag\\
&= \Theta\Bigl(\exp\bigl((\tfrac{5}{16}
  - \tfrac{\varepsilon}{2} - \tfrac{\varepsilon^2}{4})\,X\bigr)\Bigr)
  = \Theta\bigl(\exp(f(\varepsilon)\cdot X)\bigr).
\end{align}
The discriminant of $f(\varepsilon) = -\tfrac{1}{4}\varepsilon^2
- \tfrac{1}{2}\varepsilon + \tfrac{5}{16}$ yields roots at
$\varepsilon = -\tfrac{5}{2}$ and $\varepsilon = \tfrac{1}{2}$.
Since $f(0) = \tfrac{5}{16} > 0$ and $f(\tfrac{1}{2}) = 0$,
we have $f(\varepsilon) > 0$ on $(0,\tfrac{1}{2})$.
\end{proof}

\begin{remark}[The $\gamma_0$ dependence]\label{rem:gamma0}
The structural exponent $f(\varepsilon)$ depends only on $\varepsilon$
and is independent of $\mathrm{Im}(\rho_0)$. When the full
contamination-to-signal ratio is computed without simplifying assumptions,
an additional positive contribution $\gamma_0^2/4$ appears:
\[
\frac{\log R(X,\varepsilon,\gamma_0)}{X}
\;\to\; f(\varepsilon) + \frac{\gamma_0^2}{4}.
\]
Since $\gamma_0^2/4 > 0$ for all nontrivial zeros, the barrier is
\emph{strictly stronger} than the $\gamma_0 = 0$ case analyzed in
Theorem~\ref{thm:wall}. Our theoretical exponent $f(\varepsilon)$ is
therefore a lower bound on the true barrier; numerical experiments
confirm a measured exponent of approximately $0.40$ at $\varepsilon = 0.01$
compared to the theoretical $f(0.01) \approx 0.308$.
\end{remark}

\begin{observation}[Optimal Reference Scale]\label{obs:optimal_kappa}
For the single linear correction $G_X^* = F_X - c_X F_{X_0}$,
the reference scale $X_0 = \kappa X$ that minimizes the contamination
ratio satisfies $\kappa^* \to 0.9$ as $X \to \infty$.
At the optimal scale, $R(X,\varepsilon,\kappa^*)$ still grows
exponentially: scale optimization improves the constant but does not
eliminate the exponential growth.
\end{observation}

\begin{observation}[Multi-Scale Correction Failure]\label{obs:multiscale}
For a $k$-term linear correction
$G_X^{(k)}(x) = F_X(x) - \sum_{j=1}^k c_j F_{X_j}(x)$
with $\{c_j\}$ chosen to enforce $\Phi_{G^{(k)}}(1) = 0$,
numerical experiments show:
\begin{center}
\begin{tabular}{@{}cccc@{}}
\toprule
$X$ & $R^{(1)}$ & $R^{(2)}$ & Ratio \\
\midrule
20  & 518            & $2{,}110$        & 4.1$\times$ \\
50  & $6.1\times 10^6$   & $5.6\times 10^8$ & 91$\times$ \\
100 & $3.7\times 10^{13}$ & $6.2\times 10^{17}$ & $16{,}700\times$ \\
\bottomrule
\end{tabular}
\end{center}
Additional correction terms \emph{worsen} the contamination, consistent
with the interpretation that the barrier is intrinsic to the pole structure
of $\zeta(s)$ rather than an artifact of the correction method.
\end{observation}

% ============================================================
% SECTION 4: NUMERICAL EVIDENCE
% ============================================================
\section{Numerical Evidence}\label{sec:numerical}

All computations use Python 3.10 with \texttt{mpmath} at 50-digit precision.
We set $\kappa = \tfrac{1}{2}$ and $\gamma_0 = 0$ (the structurally
cleanest case; see Remark~\ref{rem:gamma0} for $\gamma_0 > 0$).

\subsection{Table 1: Contamination-to-Signal Ratio $R(X,\varepsilon)$}

\begin{table}[ht]
\centering
\caption{Contamination-to-signal ratio $R(X,\varepsilon) = \exp(f(\varepsilon)\cdot X)$
for $\kappa = 1/2$, $\gamma_0 = 0$.
Values computed from the exact closed-form $R = \exp(f(\varepsilon)\cdot X)$.}
\label{tab:R}
\begin{tabular}{@{}lcccc@{}}
\toprule
$\varepsilon \setminus X$ & 10 & 20 & 50 & 100 \\
\midrule
0.01  & 21.6          & 468            & $4.75\times 10^6$    & $2.25\times 10^{13}$ \\
0.05  & 17.8          & 317            & $1.79\times 10^6$    & $3.19\times 10^{12}$ \\
0.10  & 13.5          & 181            & $4.43\times 10^5$    & $1.96\times 10^{11}$ \\
0.20  & 6.74          & 45.4           & $1.41\times 10^4$    & $1.98\times 10^8$ \\
0.30  & 2.86          & 8.17           & $2.43\times 10^2$    & $1.20\times 10^6$ \footnote{Corrected per 50-digit verification.} \\
0.40  & 1.07          & 1.15           & 1.53                 & $2.34$ \\
0.49  & 1.0003        & 1.0006         & 1.002                & $2.11$ \\
\bottomrule
\end{tabular}
\end{table}

As $\varepsilon \to \tfrac{1}{2}$, $R \to 1$ (barrier vanishes).
For small $\varepsilon$, $R$ grows exponentially with $X$.

\subsection{Table 2: Structural Exponent $f(\varepsilon)$}

\begin{table}[ht]
\centering
\caption{Structural exponent $f(\varepsilon) = 5/16 - \varepsilon/2 - \varepsilon^2/4$ and numerical verification via $(\log R)/X$ at $X=100$.}
\label{tab:feps}
\begin{tabular}{@{}ccc@{}}
\toprule
$\varepsilon$ & $f(\varepsilon)$ (theory) & $(\log R)/X$ at $X=100$ \\
\midrule
0.01 & 0.3075 & 0.308 \\
0.05 & 0.2881 & 0.288 \\
0.10 & 0.2600 & 0.260 \\
0.20 & 0.1900 & 0.190 \\
0.30 & 0.1400 & 0.140 \\
0.40 & 0.0700 & 0.070 \\
0.49 & 0.0050 & 0.005 \\
0.50 & 0.0000 & 0.000 \\
\bottomrule
\end{tabular}
\end{table}

The theoretical and numerical values agree to displayed precision,
confirming the closed-form expression \eqref{eq:feps}.

% ============================================================
% SECTION 5: DISCUSSION
% ============================================================
\section{Discussion}\label{sec:discussion}

\subsection{Relation to Mollifier Theory}

The Levinson--Conrey method uses mollifiers of length $T^\theta$ to detect
zeros on the critical line. The barrier at $\theta = \tfrac{1}{2}$ arises
because, for $\theta < \tfrac{1}{2}$, diagonal terms in the mean-value
estimates dominate the off-diagonal contributions. Levinson \cite{Levinson1974}
obtained a proportion $> \tfrac{1}{3}$; Conrey \cite{Conrey1989} extended this
to $> \tfrac{2}{5}$ using $\theta = \tfrac{4}{7}$; Bettin and Gonek
\cite{BettinGonek2016} showed that $\theta \to \infty$ would be equivalent
to RH.

Under the correspondence $X \leftrightarrow \theta \cdot \log T$:
\begin{itemize}
\item Finite $X$ (finite $\theta$): linear correction fails
  (Theorem~\ref{thm:wall}).
\item $X \to \infty$ ($\theta \to \infty$): conditions (C1)+(C2) could in
  principle be satisfied, recovering RH (Theorem~\ref{thm:conditional}).
\end{itemize}

Both barriers share a common root: the pole of $\zeta(s)$ at $s=1$
produces a dominant contribution in the Mellin (resp.\ Fourier) domain
that any cancellation mechanism must overcome. The Gaussian barrier
$f(\varepsilon) = 0$ at $\varepsilon = \tfrac{1}{2}$ and the mollifier
barrier at $\theta = \tfrac{1}{2}$ may reflect the same Fourier-uncertainty
constraint, though establishing a formal connection remains open.

\subsection{Relation to Zero-Density Estimates}

Zero-density estimates bound $N(\sigma, T) = \#\{\rho : \beta \geq \sigma,\,
|\gamma| \leq T\}$. The classical Ingham bound $N(\sigma,T) \ll
T^{3(1-\sigma)/(2-\sigma)} \log^5 T$ and recent improvements by Guth and
Maynard \cite{GuthMaynard2024} provide statistical control over the
distribution of zeros.

Our approach is complementary: where density estimates count zeros
collectively, the pole-zero exponent inequality characterizes the
\emph{per-zero} obstruction to detection via pole-cancelled test functions.
The two perspectives address different aspects of the same underlying
difficulty.

\subsection{Limitations}

Several limitations should be noted:
\begin{enumerate}[(i)]
\item Theorem~\ref{thm:conditional} is conditional: whether test functions
  satisfying (C1)+(C2) exist is unresolved.
\item The Gaussian family is one class among many; we do not claim
  universality of the exponent $f(\varepsilon)$ across all test-function
  classes.
\item Numerical tables use $\gamma_0 = 0$ for clarity; while
  Remark~\ref{rem:gamma0} shows $\gamma_0 > 0$ only strengthens the
  barrier, a full analysis with $\gamma_0 = 14.134\ldots$ introduces
  oscillatory effects at small $X$ that warrant separate study.
\end{enumerate}

\subsection{Open Questions}

We highlight four directions for further investigation:

\begin{enumerate}[(1)]
\item \textbf{Non-Gaussian test functions.}
  Can Beurling--Selberg extremal functions or other optimized classes
  achieve $f(\varepsilon) < \tfrac{5}{16} - \tfrac{\varepsilon}{2}
  - \tfrac{\varepsilon^2}{4}$?

\item \textbf{Uniform bounds for parametric families.}
  Establishing a uniform $O(X)$ bound for $X$-dependent test-function
  families remains an open analytic problem. Such a uniform bound would
  be necessary to extend the conditional reduction in
  Theorem~\ref{thm:conditional} to parametrized families.

\item \textbf{Higher-order pole cancellation.}
  Does nonlinear correction (e.g., multiplicative or operator-theoretic)
  alter the structural exponent?

\item \textbf{Multi-zero interference.}
  Our analysis considers a single hypothetical off-critical zero.
  Interference among multiple off-critical zeros could modify the
  contamination landscape.
\end{enumerate}

\subsection{Devil's Advocate}

We address three potential objections a referee might raise:

\paragraph{``Theorem~\ref{thm:wall} is a straightforward computation.''}
The computation is elementary, but the interpretation is not.
Identifying the structural exponent $f(\varepsilon)$, its exact
vanishing at $\varepsilon^* = \tfrac{1}{2}$, and the structural
analogy with the mollifier barrier provides a new lens on a
160-year-old problem.

\paragraph{``The $O(X)$ bound in Lemma~\ref{lem:prime_bound} is weak.''}
It suffices for the contradiction argument, since the competing
term is $\Omega(e^{\varepsilon X})$. A tighter bound is desirable
but not necessary for our results.

\paragraph{``Why only Gaussians?''}
Gaussians are the canonical choice: they yield closed-form Mellin
transforms, are standard in analytic number theory (cf.\ approximate
functional equation weights), and provide a natural first benchmark.
That the barrier already appears in this best-case scenario is
informative.

% ============================================================
% APPENDICES
% ============================================================
\appendix

\section{Notation}\label{app:notation}

\begin{center}
\begin{tabular}{@{}cl@{}}
\toprule
Symbol & Meaning \\
\midrule
$\rho = \beta + i\gamma$ & Nontrivial zero of $\zeta(s)$ \\
$\varepsilon = \beta - \tfrac{1}{2}$ & Off-critical distance \\
$X$ & Scale parameter of Gaussian test function \\
$X_0 = \kappa X$ & Reference scale ($\kappa \in (0,1)$) \\
$F_X$ & Gaussian test function \eqref{eq:FX} \\
$G_X^*$ & Pole-cancelled test function \eqref{eq:correction} \\
$\Phi_X(s)$ & Mellin transform of $F_X$ \\
$c_X$ & Pole-cancellation coefficient $\Phi_X(1)/\Phi_{X_0}(1)$ \\
$R(X,\varepsilon)$ & Contamination-to-signal ratio \\
$f(\varepsilon)$ & Structural exponent $\tfrac{5}{16} - \tfrac{\varepsilon}{2}
  - \tfrac{\varepsilon^2}{4}$ \\
$C_h$ & Schwartz-norm constant for test function $h$ \\
$\Lambda(n)$ & Von Mangoldt function \\
$\psi(x)$ & Chebyshev function $\sum_{n \leq x} \Lambda(n)$ \\
$W_\infty(h)$ & Archimedean contribution in \eqref{eq:GW} \\
\bottomrule
\end{tabular}
\end{center}

\section{Proof Details}\label{app:proofs}

\subsection{Proof of Theorem~\ref{thm:conditional}}

Assume (C1) and (C2) hold, and suppose for contradiction that
$\rho_0 = \tfrac{1}{2} + \varepsilon + i\gamma$ is a nontrivial zero
of $\zeta(s)$. By the explicit formula \eqref{eq:GW}:
\[
\sum_\rho \Phi_h(\rho)
= \underbrace{\Phi_h(1)}_{= 0 \text{ by (C1)}} + \Phi_h(0)
  - \sum_n \frac{\Lambda(n)}{\sqrt{n}}[h(n)+h(1/n)] + W_\infty(h).
\]
The left-hand side includes the term $\Phi_h(\rho_0)$, which by (C2)
satisfies $|\Phi_h(\rho_0)| \geq C e^{\varepsilon X}$.
By Lemma~\ref{lem:prime_bound}, the right-hand side is $O(C_h X)$.
For sufficiently large $X$, $Ce^{\varepsilon X} > C_h X$, a contradiction.

\subsection{Derivation of the Vanishing Point $\varepsilon^* = 1/2$}

The structural exponent $f(\varepsilon) = \tfrac{5}{16} -
\tfrac{\varepsilon}{2} - \tfrac{\varepsilon^2}{4}$ is a downward-opening
quadratic. Setting $f(\varepsilon) = 0$:
\[
\varepsilon^2 + 2\varepsilon - \tfrac{5}{4} = 0
\quad\Longrightarrow\quad
\varepsilon = \frac{-2 + \sqrt{4 + 5}}{2} = \frac{-2+3}{2} = \frac{1}{2}.
\]
The negative root $\varepsilon = -\tfrac{5}{2}$ is physically irrelevant.
The geometric interpretation: $\varepsilon^* = \tfrac{1}{2}$ means
$\mathrm{Re}(\rho_0) = 1$, the location of the pole---the contamination
vanishes precisely when the ``zero'' coincides with the pole, which is
the only point where pole cancellation does not create a net signal deficit.

% ============================================================
% REFERENCES
% ============================================================
\begin{thebibliography}{99}

\bibitem{BettinGonek2016}
S.~Bettin and S.~M.~Gonek,
\emph{The $\theta = \infty$ conjecture implies the Riemann hypothesis},
Mathematika \textbf{62} (2016), no.~1, 1--19.

\bibitem{Bombieri2000}
E.~Bombieri,
\emph{The Riemann Hypothesis},
Clay Mathematics Institute Millennium Prize Problems (2000).

\bibitem{Connes2024}
A.~Connes,
\emph{Spectral triples and the Riemann zeta function},
arXiv:2310.18423 (2024).

\bibitem{Conrey1989}
J.~B.~Conrey,
\emph{More than two fifths of the zeros of the Riemann zeta function are on the critical line},
J.~Reine Angew.~Math.\ \textbf{399} (1989), 1--26.

\bibitem{Davenport2000}
H.~Davenport,
\emph{Multiplicative Number Theory},
3rd ed., Graduate Texts in Mathematics, vol.~74, Springer, 2000.

\bibitem{GoldsteinPintzYildirim2009}
D.~A.~Goldston, J.~Pintz, and C.~Y.~Y{\i}ld{\i}r{\i}m,
\emph{Primes in tuples I},
Ann.\ of Math.\ \textbf{170} (2009), no.~2, 819--862.

\bibitem{Guinand1948}
A.~P.~Guinand,
\emph{A summation formula in the theory of prime numbers},
Proc.\ London Math.\ Soc.\ \textbf{50} (1948), 107--119.

\bibitem{GuthMaynard2024}
L.~Guth and J.~Maynard,
\emph{New large value estimates for Dirichlet polynomials},
arXiv:2405.20552 (2024).

\bibitem{IwaniecKowalski2004}
H.~Iwaniec and E.~Kowalski,
\emph{Analytic Number Theory},
American Mathematical Society Colloquium Publications, vol.~53, AMS, 2004.

\bibitem{Levinson1974}
N.~Levinson,
\emph{More than one third of zeros of Riemann's zeta-function are on $\sigma = 1/2$},
Adv.\ Math.\ \textbf{13} (1974), 383--436.

\bibitem{Montgomery1973}
H.~L.~Montgomery,
\emph{The pair correlation of zeros of the zeta function},
Proc.\ Sympos.\ Pure Math.\ \textbf{24} (1973), 181--193.

\bibitem{Riemann1859}
B.~Riemann,
\emph{{\"U}ber die Anzahl der Primzahlen unter einer gegebenen Gr{\"o}sse},
Monatsberichte der Berliner Akademie (1859).

\bibitem{Selberg1942}
A.~Selberg,
\emph{On the zeros of Riemann's zeta-function},
Skr.\ Norske Vid.-Akad.\ Oslo I \textbf{10} (1942), 1--59.

\bibitem{Titchmarsh1986}
E.~C.~Titchmarsh,
\emph{The Theory of the Riemann Zeta-Function},
2nd ed.\ (revised by D.~R.~Heath-Brown), Oxford University Press, 1986.

\bibitem{Weil1952}
A.~Weil,
\emph{Sur les ``formules explicites'' de la th{\'e}orie des nombres premiers},
Comm.\ S{\'e}m.\ Math.\ Univ.\ Lund, Tome Supplementaire (1952), 252--265.

\end{thebibliography}

\end{document}
